% BANKS-SOURCE: pdf=74 print=64 kind=section id=6.1
\subsection{Free the longitudinal gauge bosons!}

% BANKS-SOURCE: pdf=74 print=64 kind=prose
Our first calculation will be the tree-level amplitude for annihilation of electron and
positron into spin-1 particles. There are two Feynman diagrams
(Figure~\ref{fig:6.1}), related
by the Bose statistics of the gauge bosons.

At tree level, we can short circuit much of the LSZ formula. The poles in external
lines sit at the bare Lagrangian masses, which are thus equal to the physical masses at
this order. Wave-function renormalization constants are all equal to one. The general
rule for external-line wave functions is that for incoming lines they are the coefficients
of annihilation operators, while for outgoing lines they are the coefficients of creation
operators, in the free fields which create and annihilate scattering states. Thus we have
\begin{equation*}
  \begin{aligned}
    u(p,s)&\;--\;\text{incoming fermion},\\
    \bar u(p,s)&\;--\;\text{outgoing fermion},\\
    v(p,s)&\;--\;\text{outgoing anti-fermion},\\
    \bar v(p,s)&\;--\;\text{incoming anti-fermion}.
  \end{aligned}
\end{equation*}
The amplitude for the first diagram is thus

% BANKS-SOURCE: pdf=74 print=64 kind=equation id=6.4
\begin{equation}
  \mathcal{M}=(-\ii e)^2\bar v(p_2,s_2)\gamma^\mu
    \frac{-\ii(\slashed p+m)}{p^2+m^2-\ii\epsilon}\gamma^\nu u(p_1,s_1)
  (\epsilon^*)_\mu(k_1,a_1)(\epsilon^*)_\nu(k_2,a_2).
  \label{eq:6.4}
\end{equation}
The momentum in the internal fermion line is
\begin{equation*}
  p^\mu=(k_1-p_2)^\mu.
\end{equation*}
The amplitude for the second diagram is the same expression with the two massive-
photon polarization vectors interchanged, and $p^\mu=k_2-p_2$. The two diagrams are to
be added, which shows that the Feynman rules implement Bose statistics.

We have written the \emph{invariant amplitude}, $\mathcal{M}$. The full S-matrix element for any process
is gotten by multiplying the invariant amplitude by
\begin{equation*}
  S=\prod_i\frac{1}{\sqrt{(2\pi)^{D-1} 2\omega_{\mathbf p_i}}}(2\pi)^D\delta^D
  \left(\sum_{\mathrm{in}}p_i-\sum_{\mathrm{out}}p_i\right)\mathcal{M}.
\end{equation*}
The product of energy factors comes from our non-relativistic state normalization, and
runs over all initial and final particles. The momentum-conservation delta function is
written here with the convention that all momenta have positive energy.

% BANKS-SOURCE: pdf=74 print=64 kind=figure id=6.1
\begin{figure}[htbp]
  \centering
  % Native TikZ reconstruction of Figure 6.1.
\begin{tikzpicture}[x=1cm,y=1cm,>=Stealth,baseline=(current bounding box.center)]
  \tikzset{BanksFermion/.style={line width=.55pt,postaction={decorate},
      decoration={markings,mark=at position .52 with {\arrow{Stealth[length=4.5pt]}}}},
    BanksPhoton/.style={line width=.55pt,decorate,
      decoration={snake,amplitude=4pt,segment length=11pt}}}
  % ---- first (uncrossed) diagram --------------------------------------
  \begin{scope}[xshift=-2.6cm]
    \draw[BanksFermion] (-1.5,0.8) -- (0,0.8);
    \draw[BanksFermion] (0,0.8) -- (0,-0.8);
    \draw[BanksFermion] (0,-0.8) -- (-1.5,-0.8);
    \draw[BanksPhoton] (0,0.8) -- (1.6,0.8);
    \draw[BanksPhoton] (0,-0.8) -- (1.6,-0.8);
  \end{scope}
  \node at (0,0) {$+$};
  % ---- second (crossed) diagram ---------------------------------------
  \begin{scope}[xshift=1.6cm]
    \draw[BanksFermion] (-1.5,0.8) -- (0,0.8);
    \draw[BanksFermion] (0,0.8) -- (0,-0.8);
    \draw[BanksFermion] (0,-0.8) -- (-1.5,-0.8);
    \draw[BanksPhoton] (0,0.8) -- (1.55,-0.85);
    \draw[BanksPhoton] (0,-0.8) -- (1.55,0.85);
  \end{scope}
\end{tikzpicture}

  \caption{Gauge-boson production in $e^+,e^-$ annihilation.}
  \label{fig:6.1}
\end{figure}

% BANKS-SOURCE: pdf=75 print=65 kind=prose
It is instructive to examine this amplitude for the case of longitudinally polarized
massive photons. The longitudinal polarization vector is
\begin{equation*}
  \epsilon^\mu(k,L)=\frac{1}{\mu}\left(|\mathbf{k}|,\frac{\omega_{\mathbf k}}{|\mathbf{k}|}\mathbf{k}\right).
\end{equation*}
Since it blows up at $\mu=0$, production of these particles would seem to rule out the
possibility of taking the massless limit. Furthermore, at high energies the dimensionless
polarization vector must blow up like $|\mathbf{k}|$. At the very least we seem to be faced by a
breakdown of perturbation theory in both the massless and the high-energy limit.

The observation that saves the day is that, at $|\mathbf{k}|\gg\mu$, we have
\begin{equation*}
  \epsilon^\mu(k,L)=\frac{k^\mu}{\mu}+o(\mu/|\mathbf{k}|).
\end{equation*}
Now consider the production amplitude for $\epsilon_1$ longitudinal and center-of-mass energy
much bigger than $\mu$. It is, in leading order,

% BANKS-SOURCE: pdf=75 print=65 kind=equation id=6.5
\begin{equation}
  -\ii e^2\bar v\left(
  \slashed k_1\frac{1}{\slashed k_1-\slashed p_2-m}\slashed\epsilon_2
  +\slashed\epsilon_2\frac{1}{\slashed k_2-\slashed p_2-m}\slashed k_1
  \right)u.
  \label{eq:6.5}
\end{equation}
Now use the identities
\begin{equation*}
  \begin{gathered}
  \slashed k_2-\slashed p_2-m=\slashed p_1-\slashed k_1-m,\\
  \slashed k_1=\slashed k_1-\slashed p_2-m+\slashed p_2+m,\\
  \slashed k_1=-(\slashed p_1-\slashed k_1-m)+\slashed p_1-m,\\
  (\slashed p_1-m)u=\bar v(\slashed p_2+m)=0.
  \end{gathered}
\end{equation*}
The terms coming from the two possible substitutions for $\slashed k_1$ each consist of one term
that cancels out the fermion propagator and another that vanishes because the external
spinors satisfy the Dirac equation. The terms with canceled propagators from the two
diagrams have equal magnitude and opposite sign. Thus, the divergent contribution to
the longitudinal production amplitude, in the zero-mass or high-energy limit, cancels
out exactly. This cancellation is the heart of the proof of renormalizability of Higgs
models.
